\section{Background}
\label{sec:background}

Among the classical concerns of information security---protecting content through encryption, preserving user privacy, and concealing the existence of communication itself---this work is concerned with the last. \textbf{Steganography} differs from encryption in a fundamental way: rather than making a message unreadable, it makes the message invisible by embedding it inside an ordinary carrier, such as an image or a piece of text, so that an observer has no reason to suspect anything is hidden. \textbf{Imperceptibility} is therefore the primary objective, not confidentiality per se.

Text is a particularly challenging carrier due to its low redundancy and strict semantic constraints. The classical ``Prisoners' Problem'' \cite{simmons1984prisoners} illustrates the goal: two parties, Alice and Bob, must exchange hidden information without alerting a watchful adversary.

Textual steganography methods are typically divided into \textbf{format-based} approaches, which exploit layout or structural features, and \textbf{linguistic-based} approaches, which modify linguistic form. Within the latter, early techniques such as \textbf{synonym substitution} embed bits by altering lexical choices, but suffer from low capacity and high detectability.

More formally, generative linguistic steganography operates through an \textbf{embedding function} $f_{emb}$ that establishes a cryptographic and linguistic mapping to hide secret messages within generated text. The algorithm takes as inputs a language model $\mathcal{L}$ and a secret message $m$ (represented as a uniformly distributed bitstream). At each generation step, rather than sampling freely from the language model's conditional distribution $p_{LM}$, the embedding function $f_{emb}$ \textbf{restricts or modifies} the token selection process according to the message bits being embedded. This specialized process is known as \textbf{steganographic sampling}, producing a steganographic text (stegotext) $y$. The security of the system depends on minimizing the divergence between the original language model distribution $p_{LM}$ and the modified distribution $q$ induced by $f_{emb}$ \cite{zhang2021provably}.

Traditional linguistic approaches offer limited embedding capacity and often leave statistical artifacts. Advances in deep learning and \textbf{Large Language Models (LLMs)} now enable generative methods that achieve higher text quality and more secure embedding. Evaluating such systems requires several dimensions of imperceptibility: \textbf{perceptual} (human naturalness), \textbf{statistical} (distributional similarity to natural text), and \textbf{cognitive} (semantic and contextual fidelity) \cite{ding2023context}.

LLM use in steganography also introduces new sources of detectability. \textbf{Hallucination}---where a model generates content that is factually unsupported, semantically inconsistent, or poorly grounded in the preceding context---is one such source. In ordinary text generation, hallucination is mainly a reliability concern. In steganographic generation, it becomes a security concern as well: even a fluent sentence can expose a hidden channel if it contains an unsupported claim, an abrupt topic shift, or a contradiction with the surrounding context. Hallucination therefore threatens \textbf{cognitive imperceptibility}, since generated text must remain contextually plausible and communicatively coherent, not merely surface-fluent.

Distinct from hallucination, the \textbf{PSIC effect} (Perceptual-Statistical Imperceptibility Conflict) \cite{yang2020vae} captures a different and arguably more fundamental tension in generative linguistic steganography. Text that is optimized for human fluency is not necessarily statistically indistinguishable from natural human text. For example, selecting only high-probability tokens can produce sentences with low perplexity and strong local readability, but it may also make the generated distribution overly smooth, regular, or concentrated compared with real human writing. Such text can appear natural to readers while remaining detectable to statistical steganalyzers \cite{yang2020vae}.

The PSIC effect therefore shows that \textbf{perceptual imperceptibility} and \textbf{statistical imperceptibility} are related but not identical goals. Human-written text often contains variation, irregularity, and locally unexpected word choices. By contrast, overly fluent model-generated text may lack this natural variance. Increasing the embedding rate can sometimes force the sampler to select from a broader range of tokens, which may better approximate the diversity of human text; however, this same process can also reduce readability or semantic coherence if the token choices become poorly aligned with the context. Thus, higher capacity does not automatically imply better security. Secure generative steganography requires balancing embedding rate, fluency, distributional similarity, and contextual consistency.

\subsection{Large Language Models as Text Generators}
\label{sec:background-llm-generators}

Large Language Models (LLMs) are autoregressive, generative systems based on the Transformer architecture \cite{vaswani2017attention} that approximate high-dimensional distributions over natural-language sequences \cite{kaptchuk2021meteor,radford2019language}. Given a prefix, an LLM emits a probability vector over the vocabulary; the next token is sampled from this vector and appended to the prefix, and the process repeats until a stopping criterion is met. During pre-training, billions of parameters are tuned on large text corpora so that the model's predictive distribution approximates patterns in the training data \cite{brown2020language}. As a consequence, modern LLMs can produce text with high fluency, coherence, and stylistic control \cite{bubeck2023sparks}. The learned representations capture stylistic and semantic regularities that generalize across domains, enabling applications requiring nuanced linguistic mimicry \cite{reif2022recipe}.

\subsection{Large Language Models in Generative Linguistic Steganography}
\label{sec:background-llm-stego}

Because text produced by modern LLMs can closely resemble ordinary human-authored text, these models provide effective engines for generative text steganography. Rather than editing an existing cover, a growing line of work treats the model itself as a steganographic sampler: the hidden message steers generation, so the stego text is drawn from a realistic communication distribution from the start. The public availability of high-quality models, together with the efficiency of sampling-based embedding, has made this route more practical than earlier designs.

LLMs like \textbf{GPT-2} \cite{radford2019language}, \textbf{LLaMA} \cite{touvron2023llama2}, and \textbf{Baichuan2} \cite{yang2023baichuan2} are commonly used as basic generative models for steganography. Most such schemes couple the language model with a steganographic code that ties message bits to token choices: at each generation step, the bits to be hidden determine which candidate token is emitted from the model's next-token distribution. However, traditional ``white-box'' methods necessitate sharing the exact language model and training vocabulary, which limits fluency, logic, and diversity compared to texts generated by modern black-box LLM APIs. These methods inevitably alter sampling probability distributions when embedding messages, posing security risks \cite{wu2024generative}. \textbf{Black-box approaches} avoid these limitations by leveraging state-of-the-art hosted models to deliver better text quality, faster generation speeds, and minimal local resource requirements. However, this paradigm trades fine-grained control over internal sampling for these practical advantages. Conversely, \textbf{white-box access} enables precise parameter control and stronger security guarantees, but demands substantial computational resources, engineering effort, and higher latency---raising deployment barriers. This trade-off between access paradigms is central to evaluating the robustness and applicability of modern linguistic steganography, and it directly motivates the black-box design adopted in Section~\ref{sec:methodology}.

New approaches, such as \textbf{LLM-Stega} \cite{wu2024generative}, explore black-box generative text steganography using the user interfaces (UIs) of LLMs, which removes any need to observe internal sampling distributions. Secret bits are encrypted and mapped onto a prearranged keyword set that steers generation, and candidate outputs from which the message cannot be recovered are discarded via reject sampling, preserving both extraction accuracy and semantic quality \cite{wu2024generative}.

Another framework, \textbf{Co-Stega}, leverages LLMs to address the challenge of low capacity in social media. It expands the space available for hiding messages through context retrieval and uses purpose-built prompts to raise the entropy of the generated replies, which in turn enlarges the embedding capacity while keeping the output fluent and topically relevant \cite{liao2024co}.

\textbf{Zero-shot linguistic steganography} instead relies on in-context learning: covertext samples are supplied in the prompt, and the stego text is produced within a question--answer (QA) exchange, so the output reads as a natural reply rather than as free-standing generated text \cite{lin2024zero}.

The increasing popularity of deep generative models has made it feasible for provably secure steganography to be applied in real-world scenarios, as they fulfill requirements for perfect samplers and explicit data distributions \cite{ding2023discop, kaptchuk2021meteor, qi2024provably}.

LLM-based steganographic methods are typically evaluated on two primary axes: imperceptibility (perceptual, statistical, and cognitive measures) and embedding capacity (e.g., bits per token or bits per word). Imperceptibility evaluations may include automatic metrics (PPL, Distinct-n, KL/JSD) as well as human judgements; embedding capacity is usually reported as bits/token or overall embedding rate. Section~\ref{sec:related_work} surveys how the literature applies, evaluates, and constrains these methods in more detail.
